% ─────────────────────────────────────────────────────────────────────────────
\chapter{About SUN-DIC}
\label{sec:about}
% ─────────────────────────────────────────────────────────────────────────────

\section{What is Digital Image Correlation?}

Digital Image Correlation (DIC) is a non-contact optical technique used to measure full-field displacement and strain on the surface of a specimen under load. A random speckle pattern is applied to the specimen surface (or may occur naturally on the surface). One or more cameras (for stereo DIC) record images as the specimen deforms.

DIC software, such as SUN-DIC, tracks small image regions called subsets between a reference image and one or more deformed images. This process, known as correlation, determines how each subset has translated and deformed during loading. The result is a dense displacement field $(u,v)$ over the selected region of interest (ROI), from which surface strains can be computed.

\section{Key Features of SUN-DIC}

SUN-DIC (Stellenbosch University Digital Image Correlation) is an open-source Python software platform for two-dimensional planar DIC analysis. It was developed at Stellenbosch University to provide a freely available, transparent, and extensible DIC solution that incorporates modern, widely used DIC algorithms.
~\\

\noindent
Key features of SUN-DIC include:

\begin{itemize}
	\item \textbf{Correlation criterion:} Zero-Mean Normalised Sum of Squared Differences (ZNSSD)
	
	\item \textbf{Optimisation algorithms:} Inverse Compositional Gauss--Newton (IC-GN) and Inverse Compositional Levenberg--Marquardt (IC-LM)
	
	\item \textbf{Shape functions:} Affine (6 DOF) and Quadratic (12 DOF)
	
	\item \textbf{Robust initialisation:} AKAZE-based initial guess strategy to improve robustness and overcome the limited convergence radius of the IC-GN and IC-LM optimisers
	
	\item \textbf{Flexible ROI definition:} Support for rectangular regions of interest, binary masks, and automatic exclusion of near-zero intensity subsets (typically corresponding to a black background)
	
	\item \textbf{Parallel execution:} Seamless parallel processing using the \code{ray} framework
	
	\item \textbf{Broad image-format support:} Compatible with all image formats supported by the \code{opencv} library
	
	\item \textbf{Multiple interfaces:} Includes both a Python API and an easy-to-use PyQt6 graphical user interface (GUI)
	
	\item \textbf{Simple installation:} Distributed through \code{pypi}

	
\end{itemize}

\section{Citation}

The software is described in detail in the following publication. Please cite this article when referencing SUN-DIC in academic publications:

\begin{grayquote}
	G.\ Venter and M.\ Neaves, ``SUN-DIC: A Python-Based Open-Source Software Tool for Digital Image Correlation,'' \textit{Advances in Engineering Software}, vol.\ 211, 2025.
	
	\href{https://doi.org/10.1016/j.advengsoft.2025.104043}{DOI: 10.1016/j.advengsoft.2025.104043}
\end{grayquote}

\section{What This Manual Covers}

This manual provides:

\begin{enumerate}
	\item Installation instructions for SUN-DIC and its dependencies
	\item A complete reference for all SUN-DIC algorithm parameters
	\item A step-by-step guide to the graphical user interface (GUI)
	\item A worked example using the Python API
	\item A concise theoretical background covering the ZNSSD criterion, subpixel interpolation, affine and quadratic shape functions, AKAZE initialisation, and the IC-GN and IC-LM optimisation algorithms
\end{enumerate}
